% One worked example, one guided example and four independent practice
% problems cover the six core families once each, in source order.
% Matching guided and independent review solutions are optional.
\makeatletter
\def\input@path{{../../../preamble/}}
\makeatother
\input{preamble.tex}
\input{graphs/highlights.tex}
\GTenHighlightsfalse % Disable instructional highlights throughout this activity.
\tikzset{gten marker/.append style={inner xsep=.6mm,inner ysep=.3mm}}
% Local C2 roles follow C1: quantities orange; probabilities in related blues.
\colorlet{cTwoMassMarker}{GTenHLProbability}
\definecolor{cTwoCDFMarker}{HTML}{43A5C9}
\newcommand{\cTwoValue}[1]{\HighlightStroke{GTenHLValue}{#1}}
\newcommand{\cTwoOutcome}[1]{\HighlightStroke{GTenHLOutcome}{#1}}
\newcommand{\cTwoMassHL}[1]{\HighlightStroke{cTwoMassMarker}{#1}}
\newcommand{\cTwoCDFHL}[1]{\HighlightStroke{cTwoCDFMarker}{#1}}
\colorlet{cTwoMassColor}{blue!75!black}
\colorlet{cTwoCDFColor}{cyan!65!blue!80!black}
\colorlet{cTwoGuideColor}{cTwoMassColor!35}
\ifGTenHighlights\else
  \colorlet{cTwoMassColor}{black}
  \colorlet{cTwoCDFColor}{black}
  \colorlet{cTwoGuideColor}{gray}
\fi
\input{graphs/G10_T2_bar_graphs.tex}
\usepackage{needspace}
\usepackage{varwidth}
\geometry{letterpaper,margin=0.25in}
\pagestyle{empty}
\linespread{1.04}
\setlength{\parskip}{4pt}
\setlength{\tabcolsep}{6pt}
\setlength{\extrarowheight}{2pt}
\setlength{\emergencystretch}{2em}
\hypersetup{pdftitle={Grade 10 C2 Classroom Activity: PMF and CDF Tables and Graphs}}
\newif\ifprintsolutions
\printsolutionsfalse % Set \printsolutionstrue to append the complete review solutions.
\newcolumntype{Y}{>{\centering\arraybackslash}X}
\newcommand{\stageheading}[1]{%
  \par%
  \Needspace{6\baselineskip}%
  \addvspace{.65\baselineskip}{\centering\large\bfseries #1%
    \par}%
  \nopagebreak[4]}
\newcommand{\problemtitle}[2]{%
  \par%
  \Needspace{4\baselineskip}%
  \addvspace{.4\baselineskip}%
  \noindent{%
    \setlength{\fboxsep}{2pt}\colorbox{black!10}{%
      \begin{varwidth}[t]{.75\linewidth}\bfseries\strut #1. #2%
      \end{varwidth}}}%
  \par%
  \nopagebreak[4]}
\newcommand{\task}[2]{%
  \par%
  \noindent\textbf{Task #1 --- \ifnum#1=1 PMF table and graph\else CDF table and graph\fi.} #2%
  \par}
\newcommand{\responseline}{%
  \par%
  \noindent%
  \rule{0pt}{1.55em}%
  \hrulefill%
  \par}
\newcommand{\worklines}[1]{\begingroup%
  \setlength{\parskip}{0pt}%
  \foreach \i in {1,...,#1}{\responseline}%
  \endgroup}
\newcommand{\blankgraph}[1]{%
  \begin{tikzpicture}%
    \begin{axis}[width=\linewidth,%
        height=4.4cm,%
        axis lines=left,%
        xmin=0,%
        xmax=10,%
        ymin=0,%
        ymax=1.05,%
        xtick=\empty,%
        ytick={0,.25,.5,.75,1},%
        grid=major,%
        grid style={gray!20},%
        xlabel={$x$ (USD)},%
        ylabel={#1},%
        tick label style={font=\small},%
        label style={font=\small}]%
    \end{axis}%
  \end{tikzpicture}}
\newcommand{\blanktable}[1]{%
  \begin{tabularx}{\linewidth}{|Y|Y|}%
    \hline #1\\%
    \hline%
    \rule{0pt}{.6cm}&\\%
    \hline%
    \rule{0pt}{.6cm}&\\%
    \hline%
    \rule{0pt}{.6cm}&\\%
    \hline%
    \rule{0pt}{.6cm}&\\%
    \hline%
    \rule{0pt}{.6cm}&\\%
    \hline%
    \rule{0pt}{.6cm}&\\%
    \hline%
  \end{tabularx}}
\pgfplotsset{classCDF/.style={width=\dimexpr\linewidth-6pt\relax,%
      height=4.15cm,%
      axis lines=left,%
      ymin=-.06,%
      ymax=1.1,%
      ytick={0,.25,.5,.75,1},%
      ylabel={$\cTwoCDFHL{F_X(x)}$},%
      xlabel={$\cTwoValue{x}$ (\cTwoValue{USD})},%
      grid=major,%
      grid style={gray!15},%
      tick label style={font=\small},%
      label style={font=\small},%
      clip=false}}
\begin{document}
{\centering{\Large\bfseries C2 -- Constructing PMF and CDF Tables and Graphs%
  \par}
Global Economics \textbar{} Grade 10 \textbar{} Term 2 \textbar{} Lesson 1%
\par}
\textbf{C2:} Interprets PMFs and CDFs by verifying their defining properties and using them to calculate
point, interval, cumulative, and complementary probabilities in financial contexts.

\vspace{-4mm}
\stageheading{Equations and basic procedure}
\noindent%
\begin{tabularx}{\linewidth}{@{}l l X@{}}%
  \toprule
  Symbol or notation & Concept & Brief meaning\\%
  \midrule
  $\cTwoValue{X} \>\> | \>\> \cTwoValue{x} \>\> | \>\> \cTwoValue{S_X}$ & \cTwoValue{Quantity} & Variable |
  Value | Support of the random variable $X$\\
  $\cTwoOutcome{\omega}$ & \cTwoOutcome{Outcome} & One complete result.\\
  $\cTwoMassHL{p_X(x) = P(X=x)}$ & \cTwoMassHL{PMF} & The \textbf{Probability Mass Function (PMF)} describes the
  probability that a discrete random variable $X$ takes a specific value $x$. \\
  $\cTwoCDFHL{F_X(x)= P(X \leq x)}$ & \cTwoCDFHL{CDF} & The \textbf{Cumulative Distribution Function (CDF)}
  describes the probability that a random variable $X$ takes a value less than or equal to a specified value
  $x$.\\%
  \bottomrule
\end{tabularx}

\task{1}{Define the random variable clearly stating what it represents and its possible values or
  \textbf{Support}. Construct a PMF table showing the probability of each possible value, and draw a labeled
  graph representing the probability distribution. Verify the PMF adds up to 1.}
\task{2}{Construct a CDF table showing the cumulative probability for each interval of values, including the
  intervals below the minimum and at or above the maximum possible values. Draw a labeled step graph
  representing the cumulative distribution, using filled and open points to indicate included and excluded
  endpoints.}

\stageheading{Worked example} \vspace{-5mm}
\problemtitle{1}{Museum kiosk payments}
\begingroup\small
\textbf{Payment fees.} \textbf{Payment fees.} A customer makes a payment at a museum kiosk using one of four methods: cash, debit
card, credit card or mobile wallet, with probabilities 0.15, 0.25, 0.35 and 0.25, respectively. The museum
pays no processing fee for cash payments, a fee of 2 USD for either type of card, and a fee of 4 USD for
mobile wallet payments. No other processing charges apply. Complete Tasks 1 and 2 to describe the
processing fee using PMF and CDF tables and graphs. Then use the CDF to find the probability that the fee
is at most 3 USD and the probability that it is at most 5 USD.


\task{1}{Let $\cTwoValue{X}$ be the \cTwoValue{processing fee} for one payment, in \cTwoValue{USD}. The
  support is $S_X=\{\cTwoValue{0},\cTwoValue{2},\cTwoValue{4}\}$.}
% Separate top-aligned blocks keep tables and graph at their native text size.
\par%
\noindent%
\begin{minipage}[t]{.28\linewidth}%
  \vspace{0pt}
  \small%
  \setlength{\tabcolsep}{3pt}
  \begin{tabularx}{\linewidth}{@{}>{\raggedright\arraybackslash}X >{\centering\arraybackslash}p{.31\linewidth}
    >{\centering\arraybackslash}p{.20\linewidth}@{}}%
    \toprule
    Complete outcome & Probability & $\cTwoValue{X}$ (\cTwoValue{USD}) \\ \midrule
    \cTwoOutcome{Cash} & $\cTwoMassHL{0.15}$ & $\cTwoValue{0}$ \\
    \cTwoOutcome{Debit card} & $\cTwoMassHL{0.25}$ & $\cTwoValue{2}$ \\
    \cTwoOutcome{Credit card} & $\cTwoMassHL{0.35}$ & $\cTwoValue{2}$ \\
    \cTwoOutcome{Wallet} & $\cTwoMassHL{0.25}$ & $\cTwoValue{4}$ \\
    \bottomrule%
  \end{tabularx}
\end{minipage}%
\hfill%
\begin{minipage}[t]{.40\linewidth}%
  \vspace{0pt}
  \small%
  \setlength{\tabcolsep}{3pt}
  \begin{tabularx}{\linewidth}{@{}>{\centering\arraybackslash}p{.15\linewidth}
    >{\centering\arraybackslash}p{.19\linewidth} >{\raggedright\arraybackslash}X@{}}%
    \toprule
    $\cTwoValue{x}$ (\cTwoValue{USD}) & $\cTwoMassHL{p_X(x)}$ & Probability calculation \\ \midrule
    $\cTwoValue{0}$ & $\cTwoMassHL{0.15}$ & $P(\text{Cash})=0.15$ \\
    $\cTwoValue{2}$ & $\cTwoMassHL{0.60}$ & $P(\text{Debit card})$\newline ${}+P(\text{Credit card})$\newline
    $=0.25+0.35=0.60$ \\
    $\cTwoValue{4}$ & $\cTwoMassHL{0.25}$ & $P(\text{Wallet})=0.25$ \\
    \bottomrule%
  \end{tabularx}
\end{minipage}%
\hfill%
\begin{minipage}[t]{.28\linewidth}%
  \vspace{0pt}\centering
  % Draw directly at this column's size so axis labels are not scaled down.
  \begin{tikzpicture}%
    \begin{axis}[
        width=\linewidth,height=4.15cm,axis lines=left,
        xmin=-1.3,xmax=5.3,ymin=0,ymax=.74,
        xtick={0,2,4},minor xtick={-1,1,3,5},ytick={0,.1,.2,.3,.4,.5,.6,.7},
        grid=both,grid style={gray!20},
        xlabel={$\cTwoValue{x}$ (\cTwoValue{USD})},ylabel={$\cTwoMassHL{p_X(x)}$},
        tick label style={font=\small},label style={font=\small},
        scaled ticks=false,clip=false]
      \draw[black,fill=cTwoMassColor!20] (axis cs:-.55,0) rectangle (axis cs:.55,.15);
      \draw[black,fill=cTwoMassColor!20] (axis cs:1.45,0) rectangle (axis cs:2.55,.60);
      \draw[black,fill=cTwoMassColor!20] (axis cs:3.45,0) rectangle (axis cs:4.55,.25);
      \addplot[only marks,mark=none,nodes near coords,
        every node near coord/.append style={font=\small},
        point meta=explicit symbolic] coordinates {(0,.15) [0.15] (2,.60) [0.60] (4,.25) [0.25]};
    \end{axis}%
  \end{tikzpicture}
\end{minipage}%
\par
\textbf{PMF check.} All masses are positive and $0.15+0.6+0.25=1$.

\task{2}{}
\par%
\noindent%
\begin{minipage}[t]{.37\linewidth}%
  \vspace{0pt}\centering
  \begin{tabular}{cc}%
    \toprule
    Cutoff interval & $\cTwoCDFHL{F_X(x)}$ \\ \midrule
    $x<0$ & $\cTwoCDFHL{0}$ \\
    $0\leq x<2$ & $\cTwoCDFHL{0.15}$ \\
    $2\leq x<4$ & $\cTwoCDFHL{0.75}$ \\
    $x\geq 4$ & $\cTwoCDFHL{1}$ \\
    \bottomrule%
  \end{tabular}
\end{minipage}%
\hfill%
\begin{minipage}[t]{.60\linewidth}%
  \vspace{0pt}\centering
  \begin{tikzpicture}%
    \begin{axis}[classCDF,%
        height=3.5cm,%
        xmin=-2,%
        xmax=6,%
        xtick={0,2,4}]
      \addplot[cTwoCDFColor,thick,<-] coordinates {(-2,0)(0,0)};
      \addplot[cTwoCDFColor,thick] coordinates {(0,0.15)(2,0.15)};
      \addplot[cTwoCDFColor,thick] coordinates {(2,0.75)(4,0.75)};
      \addplot[cTwoCDFColor,thick,->] coordinates {(4,1)(6,1)};
      \addplot[only marks,mark=*,cTwoCDFColor,mark size=1.8pt,mark options={fill=white}] coordinates
        {(0,0)(2,0.15)(4,0.75)};
      \addplot[only marks,mark=*,cTwoCDFColor,mark size=1.8pt] coordinates {(0,0.15)(2,0.75)(4,1)};
    \end{axis}%
  \end{tikzpicture}
\end{minipage}%
\par


\textbf{Fee at most 3 USD ($F_X(3)$)} Since $2\leq3<4$,
$P(X\leq3)=F_X(3)=P(X=0)+P(X=2)=0.15+0.60=0.75$.
Thus, there is a 75\% probability that the museum pays at most 3 USD in processing fees. Although $P(X=3)=0$, the CDF
includes all fees below 3 USD as well.

\textbf{Fee at most 5 USD ($F_X(5)$)} Since $5\geq4$, use the final row:
$P(X\leq5)=F_X(5)=0.15+0.60+0.25=1$.
Thus, there is a 100\% probability that the processing fee is at most 5 USD: as the maximum possible fee is 4 USD.
\clearpage
\endgroup
\vfill

\stageheading{Guided Example} \vspace{-5mm}
\problemtitle{2}{Kayak returns}
\textbf{Late fees.} A lakeside operator rents kayaks for 24 USD per hour, with a scheduled return time so
staff can prepare each kayak for the next customer. The rental charge is paid in advance; record only the
additional late fee for one completed rental. Let $t$ be the actual number of minutes after the scheduled
return time; $t\leq0$ means on time or early. Times are measured without rounding.

There is a five-minute grace period. After it, the total late fee is 6 USD for a return up to 15 minutes
late, 12 USD for more than 15 and up to 30 minutes late, and 18 USD for more than 30 and up to 45 minutes
late. Any return more than 45 minutes late costs 24 USD, the maximum late fee, even after an hour.
These are total fees, not amounts to add together. There are no other charges in this model.

For a typical rental, use the hypothetical probabilities in the outcome table. Its seven categories are
mutually exclusive and cover every possible return time.
\task{1}{Define $X$, its units and support; write the PMF table and represent it graphically. Show the outcome
  mapping and probability calculations, then verify the PMF. Group all categories that give the same fee.}

% Match the worked example: complete outcomes, grouped PMF, then graph.
% Retain the starter; students complete the fee mapping and grouped masses.
\par%
\noindent%
\begin{minipage}[t]{.28\linewidth}%
  \vspace{0pt}
  \small%
  \setlength{\tabcolsep}{3pt}
  \begin{tabularx}{\linewidth}{@{}>{\raggedright\arraybackslash}X >{\centering\arraybackslash}p{.31\linewidth}
    >{\centering\arraybackslash}p{.20\linewidth}@{}}%
    \toprule
    Return category ($t$ in min) & Probability & $\cTwoValue{X}$ (\cTwoValue{USD}) \\ \midrule
    \cTwoOutcome{On time/early: $t\leq0$} & $\cTwoMassHL{0.36}$ & $\cTwoValue{0}$ \\
    \cTwoOutcome{Grace: $0<t\leq5$} & $\cTwoMassHL{0.14}$ &  \\
    \cTwoOutcome{$5<t\leq15$} & $\cTwoMassHL{0.18}$ &  \\
    \cTwoOutcome{$15<t\leq30$} & $\cTwoMassHL{0.12}$ &  \\
    \cTwoOutcome{$30<t\leq45$} & $\cTwoMassHL{0.08}$ &  \\
    \cTwoOutcome{$45<t\leq60$} & $\cTwoMassHL{0.07}$ &  \\
    \cTwoOutcome{$t>60$} & $\cTwoMassHL{0.05}$ &  \\
    \bottomrule%
  \end{tabularx}
\end{minipage}%
\hfill%
\begin{minipage}[t]{.40\linewidth}%
  \vspace{0pt}
  \small%
  \setlength{\tabcolsep}{3pt}
  \begin{tabularx}{\linewidth}{@{}>{\centering\arraybackslash}p{.15\linewidth}
    >{\centering\arraybackslash}p{.19\linewidth} >{\raggedright\arraybackslash}X@{}}%
    \toprule
    $\cTwoValue{x}$ (\cTwoValue{USD}) & $\cTwoMassHL{p_X(x)}$ & Probability calculation \\ \midrule
    \rule{0pt}{.82cm} & & \\
    \rule{0pt}{.82cm} & & \\
    \rule{0pt}{.82cm} & & \\
    \rule{0pt}{.82cm} & & \\
    \rule{0pt}{.82cm} & & \\
    \bottomrule%
  \end{tabularx}
\end{minipage}%
\hfill%
\begin{minipage}[t]{.28\linewidth}%
  \vspace{0pt}\centering
  \begin{tikzpicture}%
    \begin{axis}[
        width=\linewidth,
        height=4.15cm,
        axis lines=left,
        xmin=-3,
        xmax=27,
        ymin=0,
        ymax=1.05,
        xtick={0,6,12,18,24},
        ytick={0,.25,.5,.75,1},
        grid=major,
        grid style={gray!20},
        xlabel={$\cTwoValue{x}$ (\cTwoValue{USD})},
        ylabel={$\cTwoMassHL{p_X(x)}$},
        tick label style={font=\small},
        label style={font=\small}]
    \end{axis}%
  \end{tikzpicture}
\end{minipage}%
\par
\textbf{PMF checks:} \hrulefill%
\par
\task{2}{Write the CDF table for every real cutoff and represent it graphically. Verify its defining
  properties and that its jumps reproduce the PMF. Interpret $F_X(12)$ in this context.}
\textbf{Prompt.} Start below the minimum at height 0; add one grouped mass at a time. Include both tails and
distinguish open and filled endpoints. Keep the horizontal axis in USD, not minutes.
\par%
\noindent%
\begin{minipage}[t]{.37\linewidth}%
  \vspace{0pt}\centering
  % Keep the removed rules' thickness as row depth to preserve table spacing.
  \begin{tabularx}{\linewidth}{|Y|Y|}%
    \hline
    Cutoff interval & $F_X(x)$ \\ \hline
    \rule{0pt}{.6cm} & \\[\arrayrulewidth]
    \rule{0pt}{.6cm} & \\[\arrayrulewidth]
    \rule{0pt}{.6cm} & \\[\arrayrulewidth]
    \rule{0pt}{.6cm} & \\[\arrayrulewidth]
    \rule{0pt}{.6cm} & \\[\arrayrulewidth]
    \rule{0pt}{.6cm} & \\ \hline
  \end{tabularx}%
\end{minipage}%
\hfill%
\begin{minipage}[t]{.60\linewidth}%
  \vspace{0pt}\centering
  \begin{tikzpicture}%
    \begin{axis}[classCDF,height=4.4cm,xmin=-3,xmax=27,xtick={0,6,12,18,24}]
    \end{axis}%
  \end{tikzpicture}
\end{minipage}%
\par
\textbf{CDF checks and contextual interpretation:}\worklines{3}
\vfill
\clearpage
\stageheading{Independent Practice}
\problemtitle{3}{Workshop scrap buyback}
\textbf{Uneven buyback payments.} A workshop scrap centre buys one standard-sized bag containing only one
material type. The exhaustive types are rejected material, paper, cardboard, plastic, steel and aluminium,
with probabilities .20, .12, .18, .15, .10, .25, respectively. Payments are 0 USD for rejected material, 4 USD
for either paper or cardboard, 9 USD for either plastic or steel, and 15 USD for aluminium. There are no other
payments. Record the buyback payment.
\task{1}{Define $X$, its units and support; write the PMF table and represent it graphically. Show the outcome
  mapping and probability calculations, then verify the PMF.}
\textbf{Calculation area:} include every complete outcome and its probability before grouping.
\worklines{4}
\par%
\noindent%
\begin{minipage}[t]{.37\linewidth}%
  \vspace{0pt}\centering
  \blanktable{$x$ (USD) & $p_X(x)$}%
\end{minipage}%
\hfill%
\begin{minipage}[t]{.60\linewidth}%
  \vspace{0pt}\centering
  \blankgraph{$p_X(x)$}
\end{minipage}%
\par
\textbf{PMF checks:} \hrulefill%
\par
\task{2}{Write the CDF table for every real cutoff and represent it graphically. Verify its defining
  properties and that its jumps reproduce the PMF. Interpret $F_X(9)$ in this context.}
\par%
\noindent%
\begin{minipage}[t]{.37\linewidth}%
  \vspace{0pt}\centering
  \blanktable{Cutoff interval & $F_X(x)$}%
\end{minipage}%
\hfill%
\begin{minipage}[t]{.60\linewidth}%
  \vspace{0pt}\centering
  \blankgraph{$F_X(x)$}
\end{minipage}%
\par
\textbf{CDF checks and contextual interpretation:}\worklines{3}
\vfill
\clearpage
\problemtitle{4}{Delivery loss claims}
\textbf{Missing probability and capped payout.} A delivery insurer can receive only 0, 1, 2, 3 or 4 claims in
one day. Their probabilities, in that order, are .15, .20, a missing probability, .20, .10. It pays 5 USD per
claim, capped at 15 USD in total for the day; it pays the smaller of 5 times the count and 15 USD. There are
no other payments. Record the daily payout.
\task{1}{Define $X$, its units and support; write the PMF table and represent it graphically. Show the outcome
  mapping and probability calculations, then verify the PMF.}
\textbf{Calculation area:} include every complete outcome and its probability before grouping.
\worklines{4}
\par%
\noindent%
\begin{minipage}[t]{.37\linewidth}%
  \vspace{0pt}\centering
  \blanktable{$x$ (USD) & $p_X(x)$}%
\end{minipage}%
\hfill%
\begin{minipage}[t]{.60\linewidth}%
  \vspace{0pt}\centering
  \blankgraph{$p_X(x)$}
\end{minipage}%
\par
\textbf{PMF checks:} \hrulefill%
\par
\task{2}{Write the CDF table for every real cutoff and represent it graphically. Verify its defining
  properties and that its jumps reproduce the PMF. Interpret $F_X(10)$ in this context.}
\par%
\noindent%
\begin{minipage}[t]{.37\linewidth}%
  \vspace{0pt}\centering
  \blanktable{Cutoff interval & $F_X(x)$}%
\end{minipage}%
\hfill%
\begin{minipage}[t]{.60\linewidth}%
  \vspace{0pt}\centering
  \blankgraph{$F_X(x)$}
\end{minipage}%
\par
\textbf{CDF checks and contextual interpretation:}\worklines{3}
\vfill
\clearpage
\problemtitle{5}{Tote bag sales}
\textbf{Independent purchases and discount.} A school enterprise offers one tote bag to each of three
visitors. Independently, each buys one with probability 0.3 or passes with probability 0.7, regardless of the
discount. Each item costs 7 USD. Only if all three buy, subtract a 5 USD discount from the combined bill.
Otherwise there is no discount. There are no other sales or charges. Record total revenue.
\task{1}{Define $X$, its units and support; write the PMF table and represent it graphically. Show the outcome
  mapping and probability calculations, then verify the PMF.}
\textbf{Calculation area:} include every complete outcome and its probability before grouping.
\par%
\noindent%
\begin{tabularx}{\linewidth}{|Y|Y|Y|Y|Y|Y|}%
  \hline Outcome & Probability & $X$ (USD) & Outcome & Probability & $X$ (USD)\\%
  \hline
  \rule{0pt}{.63cm}& & & & & \\%
  \hline
  \rule{0pt}{.63cm}& & & & & \\%
  \hline
  \rule{0pt}{.63cm}& & & & & \\%
  \hline
  \rule{0pt}{.63cm}& & & & & \\%
  \hline
\end{tabularx}
\par%
\noindent%
\begin{minipage}[t]{.37\linewidth}%
  \vspace{0pt}\centering
  \blanktable{$x$ (USD) & $p_X(x)$}%
\end{minipage}%
\hfill%
\begin{minipage}[t]{.60\linewidth}%
  \vspace{0pt}\centering
  \blankgraph{$p_X(x)$}
\end{minipage}%
\par
\textbf{PMF checks:} \hrulefill%
\par
\task{2}{Write the CDF table for every real cutoff and represent it graphically. Verify its defining
  properties and that its jumps reproduce the PMF. Interpret $F_X(14)$ in this context.}
\par%
\noindent%
\begin{minipage}[t]{.37\linewidth}%
  \vspace{0pt}\centering
  \blanktable{Cutoff interval & $F_X(x)$}%
\end{minipage}%
\hfill%
\begin{minipage}[t]{.60\linewidth}%
  \vspace{0pt}\centering
  \blankgraph{$F_X(x)$}
\end{minipage}%
\par
\textbf{CDF checks and contextual interpretation:}\worklines{3}
\vfill
\clearpage
\problemtitle{6}{Two-stage media fund}
\textbf{Signed two-stage net results.} A hypothetical classroom fund receives two cash returns. The first is
17 USD, 0 USD or 7 USD with probabilities .20, .30, .50, respectively. Independently, the second is 0 USD with
probability 0.4 or 7 USD with probability 0.6. These are all possible returns. A fixed participation cost of 7
USD is charged once. There are no other receipts or charges. Record the sum of both returns minus the cost; a
negative net result means a loss.
\task{1}{Define $X$, its units and support; write the PMF table and represent it graphically. Show the outcome
  mapping and probability calculations, then verify the PMF.}
\textbf{Calculation area:} include every complete outcome and its probability before grouping.
\par%
\noindent%
\begin{tabularx}{\linewidth}{|Y|Y|Y|Y|Y|Y|}%
  \hline Outcome & Probability & $X$ (USD) & Outcome & Probability & $X$ (USD)\\%
  \hline
  \rule{0pt}{.63cm}& & & & & \\%
  \hline
  \rule{0pt}{.63cm}& & & & & \\%
  \hline
  \rule{0pt}{.63cm}& & & & & \\%
  \hline
\end{tabularx}
\par%
\noindent%
\begin{minipage}[t]{.37\linewidth}%
  \vspace{0pt}\centering
  \blanktable{$x$ (USD) & $p_X(x)$}%
\end{minipage}%
\hfill%
\begin{minipage}[t]{.60\linewidth}%
  \vspace{0pt}\centering
  \blankgraph{$p_X(x)$}
\end{minipage}%
\par
\textbf{PMF checks:} \hrulefill%
\par
\task{2}{Write the CDF table for every real cutoff and represent it graphically. Verify its defining
  properties and that its jumps reproduce the PMF. Interpret $F_X(0)$ in this context.}
\par%
\noindent%
\begin{minipage}[t]{.37\linewidth}%
  \vspace{0pt}\centering
  \blanktable{Cutoff interval & $F_X(x)$}%
\end{minipage}%
\hfill%
\begin{minipage}[t]{.60\linewidth}%
  \vspace{0pt}\centering
  \blankgraph{$F_X(x)$}
\end{minipage}%
\par
\textbf{CDF checks and contextual interpretation:}\worklines{3}
\vfill
\clearpage
\ifprintsolutions
  \clearpage
  \stageheading{Solutions}
  \problemtitle{2}{Kayak returns}
  \begingroup\small
  \textbf{Late fees.} Use the seven return categories from the Guided Example. The prepaid rental charge
  is 24 USD per hour. The additional fee is 0 USD through the five-minute grace period, then 6, 12 and
  18 USD through 15, 30 and 45 minutes late, respectively; after 45 minutes it is capped at 24 USD.
  \task{1}{Let $\cTwoValue{X}$ be the additional late fee for one completed rental, in USD. The prepaid
    rental charge is excluded. On-time/early and grace-period returns both give 0 USD; returns more than
    45 minutes late share the 24 USD cap. Add probabilities within each of these two groups.}
  \begin{center}
    \begin{tabular}{ccc}%
      \toprule
      Return category ($t$ in min) & Probability & $\cTwoValue{X}$ (\cTwoValue{USD}) \\ \midrule
    \cTwoOutcome{On time/early: $t\leq0$} & $\cTwoMassHL{0.36}$ & $\cTwoValue{0}$ \\
    \cTwoOutcome{Grace: $0<t\leq5$} & $\cTwoMassHL{0.14}$ & $\cTwoValue{0}$ \\
    \cTwoOutcome{$5<t\leq15$} & $\cTwoMassHL{0.18}$ & $\cTwoValue{6}$ \\
    \cTwoOutcome{$15<t\leq30$} & $\cTwoMassHL{0.12}$ & $\cTwoValue{12}$ \\
    \cTwoOutcome{$30<t\leq45$} & $\cTwoMassHL{0.08}$ & $\cTwoValue{18}$ \\
    \cTwoOutcome{$45<t\leq60$} & $\cTwoMassHL{0.07}$ & $\cTwoValue{24}$ \\
    \cTwoOutcome{$t>60$} & $\cTwoMassHL{0.05}$ & $\cTwoValue{24}$ \\
      \bottomrule%
    \end{tabular}
  \end{center}
  $\cTwoValue{S_X}=\cTwoValue{\{0,6,12,18,24\}}$; $p_X(x)=0$ elsewhere. The grouped PMF is:
  \par%
  \noindent%
  \begin{minipage}[t]{.37\linewidth}%
    \vspace{0pt}\centering
    \begin{tabular}{ccc}%
      \toprule
      $\cTwoValue{x}$ (USD) & Calculation & $\cTwoMassHL{p_X(x)}$ \\ \midrule
      $\cTwoValue{0}$ & $0.36+0.14$ & $\cTwoMassHL{0.50}$ \\
      $\cTwoValue{6}$ & $0.18$ & $\cTwoMassHL{0.18}$ \\
      $\cTwoValue{12}$ & $0.12$ & $\cTwoMassHL{0.12}$ \\
      $\cTwoValue{18}$ & $0.08$ & $\cTwoMassHL{0.08}$ \\
      $\cTwoValue{24}$ & $0.07+0.05$ & $\cTwoMassHL{0.12}$ \\
      \bottomrule%
    \end{tabular}
  \end{minipage}%
  \hfill%
  \begin{minipage}[t]{.60\linewidth}%
    \vspace{0pt}\centering
    \GTenPMFGraph[color=cTwoMassColor!20,%
      xhighlight=GTenHLValue,%
      yhighlight=cTwoMassMarker,%
      width=\linewidth,%
      height=4.15cm,%
      xlabel={$\cTwoValue{x}$ (\cTwoValue{USD})},%
      ylabel={$\cTwoMassHL{p_X(x)}$},%
      xstep=6]{0,6,12,18,24}{0.50,0.18,0.12,0.08,0.12}
  \end{minipage}%
  \par
  \textbf{PMF check.} The seven outcome probabilities sum to
  $0.36+0.14+0.18+0.12+0.08+0.07+0.05=1$.
  All grouped masses are nonnegative and $0.50+0.18+0.12+0.08+0.12=1$; gaps carry no mass.
  \task{2}{Accumulate in increasing fee order: $0.50;\ 0.68;\ 0.80;\ 0.88;\ 1$.
    For example, $F_X(12)=0.50+0.18+0.12=0.80$. Each new height adds the next grouped PMF mass.}
  \par%
  \noindent%
  \begin{minipage}[t]{.37\linewidth}%
    \vspace{0pt}\centering
    \begin{tabular}{cc}%
      \toprule
      Cutoff interval & $\cTwoCDFHL{F_X(x)}$ \\ \midrule
      $x<0$ & $\cTwoCDFHL{0}$ \\
      $0\leq x<6$ & $\cTwoCDFHL{0.50}$ \\
      $6\leq x<12$ & $\cTwoCDFHL{0.68}$ \\
      $12\leq x<18$ & $\cTwoCDFHL{0.80}$ \\
      $18\leq x<24$ & $\cTwoCDFHL{0.88}$ \\
      $x\geq24$ & $\cTwoCDFHL{1}$ \\
      \bottomrule%
    \end{tabular}
  \end{minipage}%
  \hfill%
  \begin{minipage}[t]{.60\linewidth}%
    \vspace{0pt}\centering
    \begin{tikzpicture}%
      \begin{axis}[classCDF,xmin=-3,xmax=27,xtick={0,6,12,18,24}]
        \addplot[cTwoCDFColor,thick,<-] coordinates {(-3,0)(0,0)};
        \addplot[cTwoCDFColor,thick] coordinates {(0,0.50)(6,0.50)};
        \addplot[cTwoCDFColor,thick] coordinates {(6,0.68)(12,0.68)};
        \addplot[cTwoCDFColor,thick] coordinates {(12,0.80)(18,0.80)};
        \addplot[cTwoCDFColor,thick] coordinates {(18,0.88)(24,0.88)};
        \addplot[cTwoCDFColor,thick,->] coordinates {(24,1)(27,1)};
        \addplot[only marks,mark=*,cTwoCDFColor,mark size=1.8pt,mark options={fill=white}] coordinates
          {(0,0)(6,0.50)(12,0.68)(18,0.80)(24,0.88)};
        \addplot[only marks,mark=*,cTwoCDFColor,mark size=1.8pt] coordinates
          {(0,0.50)(6,0.68)(12,0.80)(18,0.88)(24,1)};
      \end{axis}%
    \end{tikzpicture}
  \end{minipage}%
  \par
  \textbf{CDF check.} Heights lie in $[0,1]$ and never decrease. Filled left/open right endpoints make the
  graph right-continuous. Horizontal tails give limits 0 and 1; the CDF reaches 1 at 24 USD. Its jumps,
  $0.50,\ 0.18,\ 0.12,\ 0.08,\ 0.12$, reproduce the PMF.
  \textbf{Interpretation.} $F_X(12)=0.80$: the probability that the additional late fee is at most 12 USD
  (a return no more than 30 minutes late, including on-time/early returns) is 0.80.
  \endgroup
  \vfill
  \clearpage
  \problemtitle{3}{Workshop scrap buyback}
  \textbf{Uneven buyback payments.} A workshop scrap centre buys one standard-sized bag containing only one
  material type. The exhaustive types are rejected material, paper, cardboard, plastic, steel and aluminium,
  with probabilities .20, .12, .18, .15, .10, .25, respectively. Payments are 0 USD for rejected material, 4 USD
  for either paper or cardboard, 9 USD for either plastic or steel, and 15 USD for aluminium. There are no other
  payments. Record the buyback payment.
  \task{1}{Let $\cTwoValue{X}$ be the buyback payment for one bag, in USD. Paper and cardboard share 4 USD, so
    their probabilities add to 0.3. Plastic and steel share 9 USD, giving 0.25. Keep the actual gaps 4, 5 and 6
    USD on both graphs.}
  \begin{center}
    \begin{tabular}{ccc}%
      \toprule
      Complete outcome & Probability & $\cTwoValue{X}$ (\cTwoValue{USD}) \\ \midrule
      \cTwoOutcome{Rejected} & $\cTwoMassHL{0.2}$ & $\cTwoValue{0}$ \\
      \cTwoOutcome{Paper} & $\cTwoMassHL{0.12}$ & $\cTwoValue{4}$ \\
      \cTwoOutcome{Cardboard} & $\cTwoMassHL{0.18}$ & $\cTwoValue{4}$ \\
      \cTwoOutcome{Plastic} & $\cTwoMassHL{0.15}$ & $\cTwoValue{9}$ \\
      \cTwoOutcome{Steel} & $\cTwoMassHL{0.1}$ & $\cTwoValue{9}$ \\
      \cTwoOutcome{Aluminium} & $\cTwoMassHL{0.25}$ & $\cTwoValue{15}$ \\
      \bottomrule%
    \end{tabular}
  \end{center}
  $\cTwoValue{S_X}=\cTwoValue{\{0,4,9,15\}}$; $p_X(x)=0$ elsewhere. The grouped PMF is:
  \par%
  \noindent%
  \begin{minipage}[t]{.37\linewidth}%
    \vspace{0pt}\centering
    \begin{tabular}{cc}%
      \toprule
      $\cTwoValue{x}$ (\cTwoValue{USD}) & $\cTwoMassHL{p_X(x)}$ \\ \midrule
      $\cTwoValue{0}$ & $\cTwoMassHL{0.2}$ \\
      $\cTwoValue{4}$ & $\cTwoMassHL{0.3}$ \\
      $\cTwoValue{9}$ & $\cTwoMassHL{0.25}$ \\
      $\cTwoValue{15}$ & $\cTwoMassHL{0.25}$ \\
      \bottomrule%
    \end{tabular}
  \end{minipage}%
  \hfill%
  \begin{minipage}[t]{.60\linewidth}%
    \vspace{0pt}\centering
    \GTenPMFGraph[color=cTwoMassColor!20,%
      xhighlight=GTenHLValue,%
      yhighlight=cTwoMassMarker,%
      width=\linewidth,%
      height=4.15cm,%
      xlabel={$\cTwoValue{x}$ (\cTwoValue{USD})},%
      ylabel={$\cTwoMassHL{p_X(x)}$},%
      xstep=1]{0,4,9,15}{0.2,0.3,0.25,0.25}
  \end{minipage}%
  \par
  \textbf{PMF check.} All masses are positive and $0.2+0.3+0.25+0.25=1$. The graph uses the actual numerical
  spacing; gaps carry no mass.
  \task{2}{Accumulate in increasing support order: $0.2;\ 0.5;\ 0.75;\ 1$. Each new height is the previous
    height plus the next PMF mass.}
  \par%
  \noindent%
  \begin{minipage}[t]{.37\linewidth}%
    \vspace{0pt}\centering
    \begin{tabular}{cc}%
      \toprule
      Cutoff interval & $\cTwoCDFHL{F_X(x)}$ \\ \midrule
      $x<0$ & $\cTwoCDFHL{0}$ \\
      $0\leq x<4$ & $\cTwoCDFHL{0.2}$ \\
      $4\leq x<9$ & $\cTwoCDFHL{0.5}$ \\
      $9\leq x<15$ & $\cTwoCDFHL{0.75}$ \\
      $x\geq 15$ & $\cTwoCDFHL{1}$ \\
      \bottomrule%
    \end{tabular}
  \end{minipage}%
  \hfill%
  \begin{minipage}[t]{.60\linewidth}%
    \vspace{0pt}\centering
    \begin{tikzpicture}%
      \begin{axis}[classCDF,%
          xmin=-4,%
          xmax=19,%
          xtick={0,4,9,15}]
        \addplot[cTwoCDFColor,thick,<-] coordinates {(-4,0)(0,0)};
        \addplot[cTwoCDFColor,thick] coordinates {(0,0.2)(4,0.2)};
        \addplot[cTwoCDFColor,thick] coordinates {(4,0.5)(9,0.5)};
        \addplot[cTwoCDFColor,thick] coordinates {(9,0.75)(15,0.75)};
        \addplot[cTwoCDFColor,thick,->] coordinates {(15,1)(19,1)};
        \addplot[only marks,mark=*,cTwoCDFColor,mark size=1.8pt,mark options={fill=white}] coordinates
          {(0,0)(4,0.2)(9,0.5)(15,0.75)};
        \addplot[only marks,mark=*,cTwoCDFColor,mark size=1.8pt] coordinates {(0,0.2)(4,0.5)(9,0.75)(15,1)};
      \end{axis}%
    \end{tikzpicture}
  \end{minipage}%
  \par
  \textbf{CDF check.} All heights lie in $[0,1]$ and never decrease. Filled left/open right endpoints make the
  graph right-continuous. Horizontal tails give limits 0 and 1; the CDF reaches 1 at 15 USD. The jumps, $0.2,\
    0.3,\ 0.25,\ 0.25$, reproduce the PMF.
  \textbf{Interpretation.} $F_X(9)=0.75$: the probability that the buyback payment is at most 9 USD (every type
  except aluminium) is 0.75.
  \vfill
  \clearpage
  \problemtitle{4}{Delivery loss claims}
  \textbf{Missing probability and capped payout.} A delivery insurer can receive only 0, 1, 2, 3 or 4 claims in
  one day. Their probabilities, in that order, are .15, .20, a missing probability, .20, .10. It pays 5 USD per
  claim, capped at 15 USD in total for the day; it pays the smaller of 5 times the count and 15 USD. There are
  no other payments. Record the daily payout.
  \task{1}{Let $\cTwoValue{X}$ be the total payout for one day, in USD. The missing probability is
    $1-(.15+.20+.20+.10)=.35$. Apply $X=\min(5N,15)$ to the claim count $N$: payouts are $0,5,10,15,15$. Four
    claims would pay 20 USD without the cap. Group the last two categories.}
  \begin{center}
    \begin{tabular}{ccc}%
      \toprule
      Complete outcome & Probability & $\cTwoValue{X}$ (\cTwoValue{USD}) \\ \midrule
      \cTwoOutcome{0 claims} & $\cTwoMassHL{0.15}$ & $\cTwoValue{0}$ \\
      \cTwoOutcome{1 claims} & $\cTwoMassHL{0.2}$ & $\cTwoValue{5}$ \\
      \cTwoOutcome{2 claims} & $\cTwoMassHL{0.35}$ & $\cTwoValue{10}$ \\
      \cTwoOutcome{3 claims} & $\cTwoMassHL{0.2}$ & $\cTwoValue{15}$ \\
      \cTwoOutcome{4 claims} & $\cTwoMassHL{0.1}$ & $\cTwoValue{15}$ \\
      \bottomrule%
    \end{tabular}
  \end{center}
  $\cTwoValue{S_X}=\cTwoValue{\{0,5,10,15\}}$; $p_X(x)=0$ elsewhere. The grouped PMF is:
  \par%
  \noindent%
  \begin{minipage}[t]{.37\linewidth}%
    \vspace{0pt}\centering
    \begin{tabular}{cc}%
      \toprule
      $\cTwoValue{x}$ (\cTwoValue{USD}) & $\cTwoMassHL{p_X(x)}$ \\ \midrule
      $\cTwoValue{0}$ & $\cTwoMassHL{0.15}$ \\
      $\cTwoValue{5}$ & $\cTwoMassHL{0.2}$ \\
      $\cTwoValue{10}$ & $\cTwoMassHL{0.35}$ \\
      $\cTwoValue{15}$ & $\cTwoMassHL{0.3}$ \\
      \bottomrule%
    \end{tabular}
  \end{minipage}%
  \hfill%
  \begin{minipage}[t]{.60\linewidth}%
    \vspace{0pt}\centering
    \GTenPMFGraph[color=cTwoMassColor!20,%
      xhighlight=GTenHLValue,%
      yhighlight=cTwoMassMarker,%
      width=\linewidth,%
      height=4.15cm,%
      xlabel={$\cTwoValue{x}$ (\cTwoValue{USD})},%
      ylabel={$\cTwoMassHL{p_X(x)}$},%
      xstep=1]{0,5,10,15}{0.15,0.2,0.35,0.3}
  \end{minipage}%
  \par
  \textbf{PMF check.} All masses are positive and $0.15+0.2+0.35+0.3=1$. The graph uses the actual numerical
  spacing; gaps carry no mass.
  \task{2}{Accumulate in increasing support order: $0.15;\ 0.35;\ 0.7;\ 1$. Each new height is the previous
    height plus the next PMF mass.}
  \par%
  \noindent%
  \begin{minipage}[t]{.37\linewidth}%
    \vspace{0pt}\centering
    \begin{tabular}{cc}%
      \toprule
      Cutoff interval & $\cTwoCDFHL{F_X(x)}$ \\ \midrule
      $x<0$ & $\cTwoCDFHL{0}$ \\
      $0\leq x<5$ & $\cTwoCDFHL{0.15}$ \\
      $5\leq x<10$ & $\cTwoCDFHL{0.35}$ \\
      $10\leq x<15$ & $\cTwoCDFHL{0.7}$ \\
      $x\geq 15$ & $\cTwoCDFHL{1}$ \\
      \bottomrule%
    \end{tabular}
  \end{minipage}%
  \hfill%
  \begin{minipage}[t]{.60\linewidth}%
    \vspace{0pt}\centering
    \begin{tikzpicture}%
      \begin{axis}[classCDF,%
          xmin=-5,%
          xmax=20,%
          xtick={0,5,10,15}]
        \addplot[cTwoCDFColor,thick,<-] coordinates {(-5,0)(0,0)};
        \addplot[cTwoCDFColor,thick] coordinates {(0,0.15)(5,0.15)};
        \addplot[cTwoCDFColor,thick] coordinates {(5,0.35)(10,0.35)};
        \addplot[cTwoCDFColor,thick] coordinates {(10,0.7)(15,0.7)};
        \addplot[cTwoCDFColor,thick,->] coordinates {(15,1)(20,1)};
        \addplot[only marks,mark=*,cTwoCDFColor,mark size=1.8pt,mark options={fill=white}] coordinates
          {(0,0)(5,0.15)(10,0.35)(15,0.7)};
        \addplot[only marks,mark=*,cTwoCDFColor,mark size=1.8pt] coordinates {(0,0.15)(5,0.35)(10,0.7)(15,1)};
      \end{axis}%
    \end{tikzpicture}
  \end{minipage}%
  \par
  \textbf{CDF check.} All heights lie in $[0,1]$ and never decrease. Filled left/open right endpoints make the
  graph right-continuous. Horizontal tails give limits 0 and 1; the CDF reaches 1 at 15 USD. The jumps, $0.15,\
    0.2,\ 0.35,\ 0.3$, reproduce the PMF.
  \textbf{Interpretation.} $F_X(10)=0.7$: the probability that the payout is at most 10 USD (at most two claims)
  is 0.7.
  \vfill
  \clearpage
  \problemtitle{5}{Tote bag sales}
  \textbf{Independent purchases and discount.} A school enterprise offers one tote bag to each of three
  visitors. Independently, each buys one with probability 0.3 or passes with probability 0.7, regardless of the
  discount. Each item costs 7 USD. Only if all three buy, subtract a 5 USD discount from the combined bill.
  Otherwise there is no discount. There are no other sales or charges. Record total revenue.
  \task{1}{Let $\cTwoValue{X}$ be the total revenue from the three visitors, in USD. Use B for buying and N for
    passing, in visitor order. Multiply within each complete sequence: e.g. $P(\mathrm{BNN})=0.3(0.7)^2=0.147$.
    The three one-buyer and three two-buyer outcomes each form a group. Only BBB receives the discount:
    $3(7)-5=16$ USD.}
  \begin{center}
    \begin{tabular}{ccc}%
      \toprule
      Complete outcome & Probability & $\cTwoValue{X}$ (\cTwoValue{USD}) \\ \midrule
      \cTwoOutcome{NNN} & $\cTwoMassHL{0.343}$ & $\cTwoValue{0}$ \\
      \cTwoOutcome{BNN} & $\cTwoMassHL{0.147}$ & $\cTwoValue{7}$ \\
      \cTwoOutcome{NBN} & $\cTwoMassHL{0.147}$ & $\cTwoValue{7}$ \\
      \cTwoOutcome{NNB} & $\cTwoMassHL{0.147}$ & $\cTwoValue{7}$ \\
      \cTwoOutcome{BBN} & $\cTwoMassHL{0.063}$ & $\cTwoValue{14}$ \\
      \cTwoOutcome{BNB} & $\cTwoMassHL{0.063}$ & $\cTwoValue{14}$ \\
      \cTwoOutcome{NBB} & $\cTwoMassHL{0.063}$ & $\cTwoValue{14}$ \\
      \cTwoOutcome{BBB} & $\cTwoMassHL{0.027}$ & $\cTwoValue{16}$ \\
      \bottomrule%
    \end{tabular}
  \end{center}
  $\cTwoValue{S_X}=\cTwoValue{\{0,7,14,16\}}$; $p_X(x)=0$ elsewhere. The grouped PMF is:
  \par%
  \noindent%
  \begin{minipage}[t]{.37\linewidth}%
    \vspace{0pt}\centering
    \begin{tabular}{cc}%
      \toprule
      $\cTwoValue{x}$ (\cTwoValue{USD}) & $\cTwoMassHL{p_X(x)}$ \\ \midrule
      $\cTwoValue{0}$ & $\cTwoMassHL{0.343}$ \\
      $\cTwoValue{7}$ & $\cTwoMassHL{0.441}$ \\
      $\cTwoValue{14}$ & $\cTwoMassHL{0.189}$ \\
      $\cTwoValue{16}$ & $\cTwoMassHL{0.027}$ \\
      \bottomrule%
    \end{tabular}
  \end{minipage}%
  \hfill%
  \begin{minipage}[t]{.60\linewidth}%
    \vspace{0pt}\centering
    \GTenPMFGraph[color=cTwoMassColor!20,%
      xhighlight=GTenHLValue,%
      yhighlight=cTwoMassMarker,%
      width=\linewidth,%
      height=4.15cm,%
      xlabel={$\cTwoValue{x}$ (\cTwoValue{USD})},%
      ylabel={$\cTwoMassHL{p_X(x)}$},%
      xstep=1]{0,7,14,16}{0.343,0.441,0.189,0.027}
  \end{minipage}%
  \par
  \textbf{PMF check.} All masses are positive and $0.343+0.441+0.189+0.027=1$. The graph uses the actual
  numerical spacing; gaps carry no mass.
  \task{2}{Accumulate in increasing support order: $0.343;\ 0.784;\ 0.973;\ 1$. Each new height is the previous
    height plus the next PMF mass.}
  \par%
  \noindent%
  \begin{minipage}[t]{.37\linewidth}%
    \vspace{0pt}\centering
    \begin{tabular}{cc}%
      \toprule
      Cutoff interval & $\cTwoCDFHL{F_X(x)}$ \\ \midrule
      $x<0$ & $\cTwoCDFHL{0}$ \\
      $0\leq x<7$ & $\cTwoCDFHL{0.343}$ \\
      $7\leq x<14$ & $\cTwoCDFHL{0.784}$ \\
      $14\leq x<16$ & $\cTwoCDFHL{0.973}$ \\
      $x\geq 16$ & $\cTwoCDFHL{1}$ \\
      \bottomrule%
    \end{tabular}
  \end{minipage}%
  \hfill%
  \begin{minipage}[t]{.60\linewidth}%
    \vspace{0pt}\centering
    \begin{tikzpicture}%
      \begin{axis}[classCDF,%
          xmin=-2,%
          xmax=18,%
          xtick={0,7,14,16}]
        \addplot[cTwoCDFColor,thick,<-] coordinates {(-2,0)(0,0)};
        \addplot[cTwoCDFColor,thick] coordinates {(0,0.343)(7,0.343)};
        \addplot[cTwoCDFColor,thick] coordinates {(7,0.784)(14,0.784)};
        \addplot[cTwoCDFColor,thick] coordinates {(14,0.973)(16,0.973)};
        \addplot[cTwoCDFColor,thick,->] coordinates {(16,1)(18,1)};
        \addplot[only marks,mark=*,cTwoCDFColor,mark size=1.8pt,mark options={fill=white}] coordinates
          {(0,0)(7,0.343)(14,0.784)(16,0.973)};
        \addplot[only marks,mark=*,cTwoCDFColor,mark size=1.8pt] coordinates {(0,0.343)(7,0.784)(14,0.973)(16,1)};
      \end{axis}%
    \end{tikzpicture}
  \end{minipage}%
  \par
  \textbf{CDF check.} All heights lie in $[0,1]$ and never decrease. Filled left/open right endpoints make the
  graph right-continuous. Horizontal tails give limits 0 and 1; the CDF reaches 1 at 16 USD. The jumps, $0.343,\
    0.441,\ 0.189,\ 0.027$, reproduce the PMF.
  \textbf{Interpretation.} $F_X(14)=0.973$: the probability that revenue is at most 14 USD (not all three
  visitors buy) is 0.973.
  \vfill
  \clearpage
  \problemtitle{6}{Two-stage media fund}
  \textbf{Signed two-stage net results.} A hypothetical classroom fund receives two cash returns. The first is
  17 USD, 0 USD or 7 USD with probabilities .20, .30, .50, respectively. Independently, the second is 0 USD with
  probability 0.4 or 7 USD with probability 0.6. These are all possible returns. A fixed participation cost of 7
  USD is charged once. There are no other receipts or charges. Record the sum of both returns minus the cost; a
  negative net result means a loss.
  \task{1}{Let $\cTwoValue{X}$ be the signed net result after both stages and the cost, in USD. For each ordered
    pair, multiply independent probabilities, then compute $X=R_1+R_2-7$. For example, $(17,0)$ has probability
    $.20(0.4)=0.08$ and net result $17+0-7=10$. The pairs $(0,7)$ and $(7,0)$ both break even; group their masses.
    Subtract the cost only once and sort signed results numerically.}
  \begin{center}
    \begin{tabular}{ccc}%
      \toprule
      Complete outcome & Probability & $\cTwoValue{X}$ (\cTwoValue{USD}) \\ \midrule
      \cTwoOutcome{$(17,0)$} & $\cTwoMassHL{0.08}$ & $\cTwoValue{10}$ \\
      \cTwoOutcome{$(17,7)$} & $\cTwoMassHL{0.12}$ & $\cTwoValue{17}$ \\
      \cTwoOutcome{$(0,0)$} & $\cTwoMassHL{0.12}$ & $\cTwoValue{-7}$ \\
      \cTwoOutcome{$(0,7)$} & $\cTwoMassHL{0.18}$ & $\cTwoValue{0}$ \\
      \cTwoOutcome{$(7,0)$} & $\cTwoMassHL{0.2}$ & $\cTwoValue{0}$ \\
      \cTwoOutcome{$(7,7)$} & $\cTwoMassHL{0.3}$ & $\cTwoValue{7}$ \\
      \bottomrule%
    \end{tabular}
  \end{center}
  $\cTwoValue{S_X}=\cTwoValue{\{-7,0,7,10,17\}}$; $p_X(x)=0$ elsewhere. The grouped PMF is:
  \par%
  \noindent%
  \begin{minipage}[t]{.37\linewidth}%
    \vspace{0pt}\centering
    \begin{tabular}{cc}%
      \toprule
      $\cTwoValue{x}$ (\cTwoValue{USD}) & $\cTwoMassHL{p_X(x)}$ \\ \midrule
      $\cTwoValue{-7}$ & $\cTwoMassHL{0.12}$ \\
      $\cTwoValue{0}$ & $\cTwoMassHL{0.38}$ \\
      $\cTwoValue{7}$ & $\cTwoMassHL{0.3}$ \\
      $\cTwoValue{10}$ & $\cTwoMassHL{0.08}$ \\
      $\cTwoValue{17}$ & $\cTwoMassHL{0.12}$ \\
      \bottomrule%
    \end{tabular}
  \end{minipage}%
  \hfill%
  \begin{minipage}[t]{.60\linewidth}%
    \vspace{0pt}\centering
    \GTenPMFGraph[color=cTwoMassColor!20,%
      xhighlight=GTenHLValue,%
      yhighlight=cTwoMassMarker,%
      width=\linewidth,%
      height=4.15cm,%
      xlabel={$\cTwoValue{x}$ (\cTwoValue{USD})},%
      ylabel={$\cTwoMassHL{p_X(x)}$},%
      xstep=1]{-7,0,7,10,17}{0.12,0.38,0.3,0.08,0.12}
  \end{minipage}%
  \par
  \textbf{PMF check.} All masses are positive and $0.12+0.38+0.3+0.08+0.12=1$. The graph uses the actual
  numerical spacing; gaps carry no mass.
  \task{2}{Accumulate in increasing support order: $0.12;\ 0.5;\ 0.8;\ 0.88;\ 1$. Each new height is the
    previous height plus the next PMF mass.}
  \par%
  \noindent%
  \begin{minipage}[t]{.37\linewidth}%
    \vspace{0pt}\centering
    \begin{tabular}{cc}%
      \toprule
      Cutoff interval & $\cTwoCDFHL{F_X(x)}$ \\ \midrule
      $x<-7$ & $\cTwoCDFHL{0}$ \\
      $-7\leq x<0$ & $\cTwoCDFHL{0.12}$ \\
      $0\leq x<7$ & $\cTwoCDFHL{0.5}$ \\
      $7\leq x<10$ & $\cTwoCDFHL{0.8}$ \\
      $10\leq x<17$ & $\cTwoCDFHL{0.88}$ \\
      $x\geq 17$ & $\cTwoCDFHL{1}$ \\
      \bottomrule%
    \end{tabular}
  \end{minipage}%
  \hfill%
  \begin{minipage}[t]{.60\linewidth}%
    \vspace{0pt}\centering
    \begin{tikzpicture}%
      \begin{axis}[classCDF,%
          xmin=-10,%
          xmax=20,%
          xtick={-7,0,7,10,17}]
        \addplot[cTwoCDFColor,thick,<-] coordinates {(-10,0)(-7,0)};
        \addplot[cTwoCDFColor,thick] coordinates {(-7,0.12)(0,0.12)};
        \addplot[cTwoCDFColor,thick] coordinates {(0,0.5)(7,0.5)};
        \addplot[cTwoCDFColor,thick] coordinates {(7,0.8)(10,0.8)};
        \addplot[cTwoCDFColor,thick] coordinates {(10,0.88)(17,0.88)};
        \addplot[cTwoCDFColor,thick,->] coordinates {(17,1)(20,1)};
        \addplot[only marks,mark=*,cTwoCDFColor,mark size=1.8pt,mark options={fill=white}] coordinates
          {(-7,0)(0,0.12)(7,0.5)(10,0.8)(17,0.88)};
        \addplot[only marks,mark=*,cTwoCDFColor,mark size=1.8pt] coordinates {(-7,0.12)(0,0.5)(7,0.8)(10,0.88)(17,1)};
      \end{axis}%
    \end{tikzpicture}
  \end{minipage}%
  \par
  \textbf{CDF check.} All heights lie in $[0,1]$ and never decrease. Filled left/open right endpoints make the
  graph right-continuous. Horizontal tails give limits 0 and 1; the CDF reaches 1 at 17 USD. The jumps, $0.12,\
    0.38,\ 0.3,\ 0.08,\ 0.12$, reproduce the PMF.
  \textbf{Interpretation.} $F_X(0)=0.5$: the probability that the net result is nonpositive (a loss or
  break-even) is 0.5.
  \vfill
  \clearpage
\fi
\end{document}
